%% Overleaf			
%% Software Manual and Technical Document Template	
%% 									
%% This provides an example of a software manual created in Overleaf.

\documentclass{../ol-softwaremanual}

% Packages used in this example
\usepackage{graphicx}  % for including images
\usepackage{microtype} % for typographical enhancements
\usepackage{minted}    % for code listings
\usepackage{amsmath}   % for equations and mathematics
\setminted{style=friendly,fontsize=\small}
\renewcommand{\listoflistingscaption}{List of Code Listings}
\usepackage{hyperref}  % for hyperlinks
\usepackage[a4paper,top=4.2cm,bottom=4.2cm,left=3.5cm,right=3.5cm]{geometry} % for setting page size and margins

\usepackage[english, greek]{babel}

\usepackage{subfig}




% Custom macros used in this example document
\newcommand{\doclink}[2]{\href{#1}{#2}\footnote{\url{#1}}}
\newcommand{\cs}[1]{\texttt{\textbackslash #1}}

\begin{document}
	
	
	\begin{titlepage}
		
		
		% Frontmatter data; appears on title page
		\title{\en Project Code \\}
		\version{0.1}
		\softwarelogo{\includegraphics[scale=0.4]{../CarBazaar_logo.png}}
	\end{titlepage}
	
	
	\maketitle
	
	\newpage
	
	\center{\textbf{Μέλη Ομάδας}}
	
	\vspace{20pt}
	
	
	
	\begin{table}[htbp!]
		
		\begin{tabular}{llll}
			Μεμελετζόγλου Χαρίλαος & 1069364 & \en st1069364@ceid.upatras.gr & 4o Έτος   \\ 
			\\ Λέκκας Γεώργιος      &      1067430    &   \en st1067430@ceid.upatras.gr & 4o Έτος  \\
			\\ Γιαννουλάκης Ανδρέας        &   1067387       & \en st1067387@ceid.upatras.gr & 4o Έτος           \\
			\\ Κανελλόπουλος Ιωακείμ        &  1070914        &    \en st1070914@ceid.upatras.gr & 4o Έτος        \\ 
		\end{tabular}
	\end{table}
	
	\center{\textbf{Υπεύθυνοι Παρόντος Τεχνικού Κειμένου}}
	
	\vspace{20pt}
	
	\begin{table}[htbp!]
		\begin{tabular}{ll}
			Μεμελετζόγλου Χαρίλαος & \en Editor \\
			\\ Λέκκας Γεώργιος      &   \en  Peer Reviewer \\
			\\ Γιαννουλάκης Ανδρέας & \en Peer Reviewer \\
			\\ Καννελόπουλος Ιωακείμ & \en Peer Reviewer \\ 
		\end{tabular}
	\end{table}
	
	
	
	\center{\textbf{Εργαλεία που χρησιμοποιήθηκαν}}
	
	\vspace{20pt}
	\flushleft
	Χρησιμοποιήθηκε το \en \doclink{https://www.overleaf.com/}{Overleaf} \gr και το \en \doclink{https://www.texstudio.org/}{TexStudio} \gr για την συγγραφή του \LaTeX\ κώδικα. \break
	
	Για την δημιουργία του \en UI \gr, χρησιμοποιείται το \en \textbf{PyQt} cross-platform Qt Framework \gr. \break
	
	Για την συγγραφή του \en Python \gr κώδικα, χρησιμοποιείται το \en IDE \doclink{https://www.jetbrains.com/pycharm/}{PyCharm} \gr και το \en \doclink{https://code.visualstudio.com/}{VSCode} \gr .         \\ \break
	
	Για την αποδοτικότερη δημιουργία των οθονών, χρησιμοποιείται το \en \doclink{https://doc.qt.io/qt-5/qtdesigner-manual.html}{QtDesigner} \gr, το οποίο παράγει τα \en .ui \gr, αρχεία που βρίσκονται στο \en src/ui/qt\_ui directory \gr . \break
		
	
	\newpage
	
	
	\center{\textbf{Περιγραφή κώδικα στο \en Github \gr}}
	
	\flushleft
	
	Στο \en \doclink{https://github.com/st1069364/CarBazaar}{repo} \gr του \en Project \gr, υπάρχει μια πρώτη μετατροπή των \en Mockup \gr σε πραγματικές \en UI \gr οθόνες. \break
	
	Ο κώδικας βρίσκεται στο \en \textbf{src} directory \gr . \break
	
	Συγκεκριμένα, υπάρχει μια υλοποίηση του \en User Menu \gr χρησιμοποιώντας \en PyQt \gr (αρχείο \en ui/user\_menu\_qt.py) \gr \break 
	
	Στο \en directory \textbf{ui/qt\_ui} \gr υπάρχουν τα απαραίτητα \en resources \gr, καθώς και το \en .ui \gr αρχείο που προκύπτει από το \en Qt Designer \gr. \break
	
	Επίσης, υπάρχουν υλοποιήσεις των \en User Menu, Login Screen \& Signup Screen \gr, χρησιμοποιώντας \en Tkinter \gr (στο μέλλον θα αντικατασταθούν από \en PyQt \gr εκδοχές). \break
	
	

	

	
	
	
	
	
\end{document}
